\documentclass[11pt]{article}
\usepackage[a4paper,margin=1in]{geometry}
\usepackage{amsmath,amssymb,amsthm}
\usepackage{enumitem}
\usepackage{hyperref}

\title{AxlerLean Exercises\\(Lean + LaTeX Companion)}
\author{TheMathBook / Part I / Linear Algebra}
\date{\today}

\newtheorem*{problem}{Problem}
\newtheorem*{solution}{Solution}

\begin{document}
\maketitle

\section*{How This File Matches Lean}
Source of truth is:
\begin{itemize}[leftmargin=1.5em]
  \item \texttt{AxlerLean/Exercises.lean}
\end{itemize}

This LaTeX file is a readable companion. The theorem names below match Lean theorem IDs.

\section*{Basic (p1--p6)}

\begin{problem}[p1\_zero\_smul\_vec]
Show that $0 \cdot v = 0$ coordinatewise.
\end{problem}
\begin{solution}
For each coordinate $i$, we have $0 \cdot v_i = 0$. Hence the two functions are equal by extensionality.
\end{solution}

\begin{problem}[p2\_vec\_ext]
If all coordinates agree, prove $u=v$.
\end{problem}
\begin{solution}
Assume $\forall i,\ u_i=v_i$. Equality of vectors follows from coordinatewise equality.
\end{solution}

\begin{problem}[p3\_id\_apply]
Identity linear map acts as identity.
\end{problem}
\begin{solution}
By definition of \texttt{LinearMap.id}, \texttt{toFun v = v}.
\end{solution}

\begin{problem}[p4\_eval\_at\_zero]
For $f(x)=ax+b$, show $f(0)=b$.
\end{problem}
\begin{solution}
$f(0)=a\cdot0+b=b$.
\end{solution}

\begin{problem}[p5\_add\_zero\_vec]
Show that $v+0=v$ coordinatewise.
\end{problem}
\begin{solution}
For each coordinate, $v_i+0=v_i$. Thus the vector functions are equal.
\end{solution}

\begin{problem}[p6\_det\_zero\_matrix]
In the $2\times2$ model, prove $\det\!\begin{bmatrix}0&0\\0&0\end{bmatrix}=0$.
\end{problem}
\begin{solution}
Using $\det\begin{bmatrix}a&b\\c&d\end{bmatrix}=ad-bc$, substitute $a=b=c=d=0$.
\end{solution}

\section*{Standard (p7--p12)}

\begin{problem}[p7\_eigen\_nonzero]
If $v$ is an eigenvector by definition, show $v\neq0$.
\end{problem}
\begin{solution}
Immediate from the first conjunct in the Lean predicate \texttt{IsEigenvector}.
\end{solution}

\begin{problem}[p8\_dot\_zero\_left]
For the 2D dot model, prove $\mathrm{dot}(0,u)=0$.
\end{problem}
\begin{solution}
Expand the definition and simplify both terms to zero.
\end{solution}

\begin{problem}[p9\_map\_zero]
Every linear map sends the zero vector to zero.
\end{problem}
\begin{solution}
Use linearity with scalar $0$: $T(0\cdot \mathbf{1})=0\cdot T(\mathbf{1})=0$.
\end{solution}

\begin{problem}[p10\_conj\_involutive]
Conjugation is involutive in the Pair model.
\end{problem}
\begin{solution}
Apply conjugation twice: $(\mathrm{re},\mathrm{im})\mapsto(\mathrm{re},-\mathrm{im})\mapsto(\mathrm{re},\mathrm{im})$.
\end{solution}

\begin{problem}[p11\_comp\_id\_right]
Show $f\circ \mathrm{id}=f$.
\end{problem}
\begin{solution}
Directly from the proved lemma \texttt{LinearMap.comp\_id\_right}.
\end{solution}

\begin{problem}[p12\_eval\_at\_one]
For $f(x)=ax+b$, show $f(1)=a+b$.
\end{problem}
\begin{solution}
$f(1)=a\cdot1+b=a+b$.
\end{solution}

\section*{Challenge (p13--p18)}

\begin{problem}[p13\_comp\_id\_left]
Show $\mathrm{id}\circ f=f$.
\end{problem}
\begin{solution}
Directly from \texttt{LinearMap.comp\_id\_left}.
\end{solution}

\begin{problem}[p14\_comp\_assoc]
Show associativity of composition.
\end{problem}
\begin{solution}
Use \texttt{LinearMap.comp\_assoc}.
\end{solution}

\begin{problem}[p15\_vec\_zero\_unique]
If all coordinates of $v$ are zero, show $v=0$.
\end{problem}
\begin{solution}
Apply function extensionality with the coordinate hypothesis.
\end{solution}

\begin{problem}[p16\_dot\_comm]
In the 2D model, prove $\mathrm{dot}(u,v)=\mathrm{dot}(v,u)$.
\end{problem}
\begin{solution}
Commutativity of integer multiplication on each coordinate term.
\end{solution}

\begin{problem}[p17\_det\_equal\_rows]
If two rows are equal, determinant is zero (2x2 model).
\end{problem}
\begin{solution}
Substitute $(a,b)$ and $(a,b)$ in $ad-bc$, giving $ab-ba=0$.
\end{solution}

\begin{problem}[p18\_conj\_re\_fixed]
Real part is invariant under conjugation.
\end{problem}
\begin{solution}
From the definition \texttt{conj (re,im) = (re,-im)}.
\end{solution}

\section*{Build}
From this directory:
\begin{verbatim}
pdflatex Exercises.tex
\end{verbatim}
or
\begin{verbatim}
xelatex Exercises.tex
\end{verbatim}

\end{document}
